\section{The discrete logarithm problem}

\subsection{Formal definition}

\begin{definition}[Discrete logarithm problem]
Let $\G=\langle g\rangle$ be a cyclic group of order $r$. Given $g$ and
$h\in\G$, the discrete logarithm problem is to find
\[
  x\in\Z/r\Z\quad\text{such that}\quad h=g^x.
\]
In additive notation, given $G$ and $Q\in\langle G\rangle$, find $x$ such that
\[
  Q=[x]G.
\]
\end{definition}

The group, generator, and generator order are part of the problem statement.
The answer is unique only modulo $\ord(g)$. If the target is outside the
generated subgroup, no answer exists.

Forward evaluation uses square-and-multiply or double-and-add and requires
$O(\log r)$ group operations. Reversal has no comparable generic algorithm.
This asymmetry is computational, not algebraic: the forward map is a homomorphism
from $\Z/r\Z$ onto $\langle g\rangle$.

\subsection{Generic attacks}

\begin{center}
\begin{tabular}{lll}
\toprule
Method & Time & Storage \\
\midrule
exhaustive search & $O(r)$ & $O(1)$ \\
baby-step--giant-step & $O(\sqrt r)$ & $O(\sqrt r)$ \\
Pollard rho & expected $O(\sqrt r)$ & $O(1)$ \\
Pohlig--Hellman & controlled by the largest prime factor of $r$ & varies \\
\bottomrule
\end{tabular}
\end{center}

Baby-step--giant-step writes $x=im+j$ with $m=\lceil\sqrt r\rceil$ and
looks for a collision
\[
  h(g^{-m})^i=g^j.
\]
Pollard rho replaces the stored table by a pseudorandom walk and a cycle-finding
method \cite{pollard1978rho}. Generic-group lower bounds explain why the square
root repeatedly appears \cite{shoup1997lower}.

Pohlig--Hellman reduces a DLP of order
$r=\prod_i q_i^{e_i}$ to DLPs in the prime-power factors. A secure group
therefore needs a large prime factor; a large ambient field by itself is not
enough.

Finite-field multiplicative groups also admit index-calculus-style algorithms
that exploit their representation and beat the generic square-root bound.
No analogous general subexponential method is known for a properly chosen
elliptic-curve group. This is why a 255-bit elliptic-curve field can target a
security level that would require a much larger finite-field Diffie--Hellman
modulus.

\subsection{DLP, CDH, and DDH are different problems}

Let $g$ generate a group of prime order $r$ and choose unknown
$a,b\in\Z/r\Z$.

\begin{definition}[Computational Diffie--Hellman]
The $\CDH$ problem maps
\[
  (g,g^a,g^b)\longmapsto g^{ab}.
\]
In additive notation it maps $(G,[a]G,[b]G)$ to $[ab]G$.
\end{definition}

\begin{definition}[Decisional Diffie--Hellman]
The $\DDH$ problem asks whether a supplied $T$ equals $g^{ab}$, given
$(g,g^a,g^b,T)$.
\end{definition}

A DLP solver gives a CDH solver: recover $a$, then compute $(g^b)^a$.
A CDH solver gives a DDH solver: compute $g^{ab}$ and compare it with $T$.
Thus there are solver implications
\[
  \DLP\Longrightarrow\CDH\Longrightarrow\DDH,
\]
but these problems are not known to be equivalent in arbitrary groups. In
some pairing-friendly groups, for example, DDH is easy while CDH may remain
hard. Saying ``DLP is hard'' is therefore not itself a proof of every
Diffie--Hellman assumption.

\begin{takeaway}
Curve25519 does not introduce a new kind of logarithm. It supplies a carefully
chosen cyclic subgroup in which the ordinary DLP is written additively and is
called ECDLP.
\end{takeaway}
